\documentclass[11pt,reqno]{amsart}

\usepackage[T1]{fontenc}
\usepackage{lmodern}
\usepackage{microtype}
\usepackage{mathtools}
\usepackage{amssymb}
\usepackage{enumitem}
\usepackage{xcolor}
\usepackage{hyperref}

\hypersetup{
  colorlinks=true,
  linkcolor=blue!55!black,
  citecolor=blue!55!black,
  urlcolor=blue!55!black,
  pdftitle={The monodromy-character law for joint Hurwitz hypertranscendence},
  pdfauthor={Michael Haufschild}
}

\setlist[enumerate]{leftmargin=2.2em,itemsep=0.25em,topsep=0.4em}
\setlist[itemize]{leftmargin=2.0em,itemsep=0.25em,topsep=0.4em}
\setlength{\emergencystretch}{2em}

\newtheorem{theorem}{Theorem}[section]
\newtheorem{proposition}[theorem]{Proposition}
\newtheorem{lemma}[theorem]{Lemma}
\newtheorem{corollary}[theorem]{Corollary}
\theoremstyle{definition}
\newtheorem{definition}[theorem]{Definition}
\theoremstyle{remark}
\newtheorem{remark}[theorem]{Remark}

\newcommand{\C}{\mathbb C}
\newcommand{\Q}{\mathbb Q}
\newcommand{\N}{\mathbb N}
\newcommand{\Z}{\mathbb Z}
\newcommand{\HH}{\mathbb H}
\newcommand{\Ainf}{\mathcal A_{\infty}}
\newcommand{\Minf}{\mathcal M_{\infty}}
\newcommand{\mc}{\mathrm{MC}}
\newcommand{\dd}{\,\mathrm d}
\DeclareMathOperator{\Frac}{Frac}
\DeclareMathOperator{\trdeg}{trdeg}

\title[Joint Hurwitz hypertranscendence]
  {The monodromy-character law for joint Hurwitz hypertranscendence}

\author{Michael Haufschild}
\date{}

\subjclass[2020]{Primary 11M35, 12H05; Secondary 30F99, 39A06}
\keywords{Hurwitz zeta function, hypertranscendence, differential algebraic independence,
monodromy, Zariski density, Vandermonde matrix}

\begin{document}

\begin{abstract}
Let $I$ be a finite set, let $(s_j)_{j\in I}$ be complex numbers with
$\operatorname{Re}s_j>1$, and put $q_j=\exp(-2\pi i s_j)$.  Assume that the
monomial-character map
\[
  \N^I\longrightarrow\C^\times,
  \qquad \alpha\longmapsto q^\alpha:=\prod_{j\in I}q_j^{\alpha_j},
\]
is injective.  We prove that the germs of the functions
$a\mapsto\zeta(s_j,a)$ are differentially algebraically independent over
$\C(a)$.  Equivalently, for every $N\geq 0$, all functions
\[
  \zeta(s_j+k,a),\qquad j\in I,\quad 0\leq k\leq N,
\]
are algebraically independent over $\C(a)$.  Rational linear independence of
$1$ and the exponents $(s_j)_{j\in I}$ implies the monomial-character
hypothesis.  The result extends, in the finite-subfamily sense, to an
arbitrarily indexed rationally independent exponent family.  The proof uses
simultaneous monodromy at the branch points $0,-1,-2,\ldots$, Zariski density
of a common cyclic power orbit, independent branch-point columns, and a
block-Vandermonde automorphism.
\end{abstract}

\maketitle

\section{Introduction}

For $\operatorname{Re}s>1$ and $\operatorname{Im}a>0$, the Hurwitz zeta
function in the shift variable is represented by the normally convergent
series
\begin{equation}\label{eq:hurwitz-series}
  \zeta(s,a)=\sum_{n=0}^{\infty}(n+a)^{-s},
\end{equation}
where powers are defined using the principal logarithm.  The dependence on
the two variables is asymmetric.  In the Dirichlet-series variable $s$, the
terms have fixed frequencies $\log(n+a)$.  In the shift variable $a$, the
function has an infinite arithmetic array of branch points at
$0,-1,-2,\ldots$, and differentiation raises the exponent:
\begin{equation}\label{eq:derivative-law-intro}
  \partial_a^k\zeta(s,a)=(-1)^k(s)_k\zeta(s+k,a),
  \qquad
  (s)_k=s(s+1)\cdots(s+k-1).
\end{equation}
The first derivative is classical; see, for example,
\cite[Eq.~25.11.17]{DLMF}.  Formula \eqref{eq:derivative-law-intro} follows by
iteration.

The problem considered here is joint and nonlinear.  Fix several exponents
$s_j$.  Can the functions $a\mapsto\zeta(s_j,a)$ satisfy a mixed algebraic
differential equation over $\C(a)$?  Separate hypertranscendence of the
individual functions would not answer this question: it excludes an equation
in one differential indeterminate at a time, but not an equation coupling
several exponent families.

The branch geometry identifies the relevant arithmetic datum.  A positive
loop around $a=-n$ multiplies the singular germ contributed by the $n$th
summand in the $j$th family by
\begin{equation}\label{eq:qj-intro}
  q_j=\exp(-2\pi i s_j).
\end{equation}
The same geometric loop acts simultaneously on every $j$.  Thus one does not
obtain independent powers of the $q_j$ at a fixed branch point.  The native
orbit is the common-time orbit
\[
  m\longmapsto(q_j^m)_{j\in I},\qquad m\in\N.
\]
Its polynomial Zariski density is equivalent to separation of its monomial
characters.  Different branch points, in contrast, can be looped around
independently.  Their product orbit supplies one common-time column per branch
point.  A block-Vandermonde map then converts the branch parameters into
arbitrary finite ladder values.

This mechanism yields the following law.  Monomial separation of the common
monodromy character forces complete joint differential freeness of the
corresponding Hurwitz family.  Rational linear independence of
$1,(s_j)_{j\in I}$ is a directly checkable sufficient condition.  The theorem
holds over $\C(a)$, not merely over the constant field, and it includes all
mixed derivatives of all finite orders.

The analytic continuation and linear monodromy spaces of the Hurwitz-Lerch
family were developed systematically by Lagarias and Li
\cite{LagariasLiII,LagariasLiIII}.  Related monodromy calculations for
fractional polylogarithms appear in Costin and Garoufalidis
\cite{CostinGaroufalidis}.  Joyner proved algebraic independence over $\C$ for
an integral-index Hurwitz polyzeta algebra by difference equations and
stuffle methods \cite{Joyner}.  Differential-transcendence results of Amou
and Katsurada concern generalized Dirichlet series in the variable $s$
\cite{AmouKatsurada}; universality-based functional-independence results also
vary the Dirichlet-series variable \cite{Laurincikas}.  The general
shift-equation theorem of Adamczewski, Dreyfus, and Hardouin concerns
homogeneous linear shift equations over the rational function field
\cite{ADH}.  For nonintegral fixed $s$, the elementary Hurwitz relation
\[
  \zeta(s,a+1)-\zeta(s,a)=-a^{-s}
\]
has a forcing term outside $\C(a)$.  The present result instead uses the full
branch array in the shift variable and gives a joint monodromy-character
criterion.

The paper is organized as follows.  Section~\ref{sec:field} fixes the
differential-field setting.  Section~\ref{sec:statements} states the main
results.  Section~\ref{sec:orbits} proves the power-orbit lemmas.
Section~\ref{sec:deformation} establishes simultaneous branch deformation.
Section~\ref{sec:main-proof} proves ladder independence, and
Section~\ref{sec:differential} derives the differential-field and
arbitrary-family forms.  The final section records the exact boundary of the
method.

\section{The upper-end differential field}\label{sec:field}

Let
\[
  \HH=\{a\in\C:\operatorname{Im}a>0\}.
\]
For $H\in\mathbb R$, write
\[
  \HH_H=\{a\in\C:\operatorname{Im}a>H\}.
\]

\begin{definition}[Analytic germs at the upper end]
Two functions analytic on upper half-planes are equivalent if they agree on
some upper half-plane on which both are defined.  Let $\Ainf$ denote the ring
of the resulting germs.  Its fraction field is denoted by
\[
  \Minf=\Frac(\Ainf).
\]
\end{definition}

The identity theorem shows that $\Ainf$ is an integral domain.  Indeed, if
the product of two analytic upper-end germs is zero, choose one half-plane on
which both representatives are analytic and their product vanishes.  On the
connected domain $\HH_H$, one of the two analytic representatives vanishes
identically.  Hence one of the two germs is zero.

Ordinary complex differentiation descends to $\Ainf$: if two analytic
representatives agree on an upper half-plane, their derivatives agree there
as well.  The quotient rule extends this derivation to $\Minf$.  Polynomial
evaluation embeds $\C[a]$ faithfully in $\Ainf$, because a nonzero polynomial
cannot vanish on an upper half-plane.  Passing to fraction fields gives a
faithful differential-field embedding
\begin{equation}\label{eq:ratfunc-embedding}
  \C(a)\hookrightarrow\Minf.
\end{equation}

For $\operatorname{Re}s>1$, the function $a\mapsto\zeta(s,a)$ is analytic on
$\HH$, hence defines an element of $\Ainf$ and therefore of $\Minf$.  We use
the same notation $\zeta(s,a)$ for this germ.

\begin{definition}[Differential algebraic independence]
Let $(K,\partial)$ be a differential field extension of a field $F$.  A
family $(x_j)_{j\in I}\subset K$ is \emph{differentially algebraically
independent over $F$} if the family of all jets
\[
  \{\partial^k x_j:j\in I,\ k\in\N\}
\]
is algebraically independent over $F$ in the finite-subfamily sense.
\end{definition}

\section{Main results}\label{sec:statements}

Let $I$ be a finite set.  For $s=(s_j)_{j\in I}$ with
$\operatorname{Re}s_j>1$, define
\begin{equation}\label{eq:monodromy-vector}
  q_j=\exp(-2\pi i s_j),\qquad j\in I.
\end{equation}
For $\alpha=(\alpha_j)_{j\in I}\in\N^I$, put
\[
  q^\alpha=\prod_{j\in I}q_j^{\alpha_j}.
\]

\begin{definition}[Monodromy-character separation]\label{def:MC}
The exponent family $s=(s_j)_{j\in I}$ satisfies condition $\mc$ if
\begin{equation}\label{eq:MC}
  \alpha\longmapsto q^\alpha
  \quad\text{is injective on }\N^I.
\end{equation}
\end{definition}

For $N\geq0$, write
\begin{equation}\label{eq:ladder}
  \mathcal L_N(s)=
  \bigl(\zeta(s_j+k,a)\bigr)_{j\in I,\ 0\leq k\leq N}.
\end{equation}

\begin{theorem}[Monodromy-character law]\label{thm:main}
Let $I$ be finite, let $s_j\in\C$ satisfy
$\operatorname{Re}s_j>1$, and assume condition $\mc$.  Then the elements
\[
  \zeta(s_j,a)\in\Minf,\qquad j\in I,
\]
are differentially algebraically independent over $\C(a)$.  Equivalently,
for every $N\geq0$, the entries of $\mathcal L_N(s)$ are algebraically
independent over $\C(a)$.
\end{theorem}

\begin{theorem}[Polynomial-coefficient form]\label{thm:polynomial}
Under the hypotheses of Theorem~\ref{thm:main}, fix $N\geq0$.  If
\[
  P\in\C[a][Y_{j,k}:j\in I,\ 0\leq k\leq N]
\]
satisfies
\begin{equation}\label{eq:polynomial-relation}
  P\bigl(a;(\zeta(s_j+k,a))_{j,k}\bigr)=0
  \qquad(a\in\HH),
\end{equation}
then $P=0$.
\end{theorem}

\begin{corollary}[Arithmetic form]\label{cor:arithmetic}
Let $I$ be finite and let $s_j\in\C$ satisfy
$\operatorname{Re}s_j>1$.  If the family consisting of $1$ and
$(s_j)_{j\in I}$ is linearly independent over $\Q$, then the conclusions of
Theorems~\ref{thm:main} and~\ref{thm:polynomial} hold.
\end{corollary}

\begin{corollary}[Arbitrary exponent families]\label{cor:arbitrary}
Let $K$ be any index set, and let $(s_j)_{j\in K}$ satisfy
$\operatorname{Re}s_j>1$.  Assume that $1$ together with all $s_j$ is
linearly independent over $\Q$.  Then
\[
  \{\partial_a^k\zeta(s_j,a):j\in K,\ k\in\N\}
\]
is algebraically independent over $\C(a)$ in the finite-subfamily sense.
The same conclusion holds for
\[
  \{\zeta(s_j+k,a):j\in K,\ k\in\N\}.
\]
\end{corollary}

\section{Cyclic power orbits}\label{sec:orbits}

The common loop count at one branch point is the central constraint.  The
following lemma converts character separation into Zariski density.

\begin{lemma}[Power-orbit density]\label{lem:power-orbit}
Let $I$ be finite and let $q=(q_j)_{j\in I}\in(\C^\times)^I$.  Suppose that
$\alpha\mapsto q^\alpha$ is injective on $\N^I$.  If
$R\in\C[X_j:j\in I]$ satisfies
\begin{equation}\label{eq:orbit-vanish}
  R((q_j^m)_{j\in I})=0\qquad(m\in\N),
\end{equation}
then $R=0$.
\end{lemma}

\begin{proof}
Write
\[
  R=\sum_{\alpha\in A}c_\alpha X^\alpha,
\]
where $A\subset\N^I$ is finite.  Equation
\eqref{eq:orbit-vanish} becomes
\begin{equation}\label{eq:power-sum}
  \sum_{\alpha\in A}c_\alpha(q^\alpha)^m=0
  \qquad(m\in\N).
\end{equation}
The numbers $q^\alpha$, $\alpha\in A$, are pairwise distinct.  Evaluating
\eqref{eq:power-sum} for $m=0,1,\ldots,|A|-1$ gives a square Vandermonde
system with determinant
\[
  \prod_{\substack{\alpha,\beta\in A\\ \alpha<\beta}}
  (q^\beta-q^\alpha)\neq0,
\]
after choosing any ordering of $A$.  Hence every $c_\alpha$ is zero.
\end{proof}

\begin{lemma}[Product-grid density]\label{lem:product-grid}
Under the hypotheses of Lemma~\ref{lem:power-orbit}, let $L\geq1$.  Suppose
that
\[
  R\in\C[X_{j,n}:j\in I,\ 0\leq n<L]
\]
vanishes at every point
\begin{equation}\label{eq:product-grid}
  X_{j,n}=q_j^{m_n},
  \qquad (m_0,\ldots,m_{L-1})\in\N^L.
\end{equation}
Then $R=0$.
\end{lemma}

\begin{proof}
Fix orbit values in columns $1,\ldots,L-1$ and regard $R$ as a polynomial
in column $0$.  It vanishes on the full cyclic power orbit, so
Lemma~\ref{lem:power-orbit} makes it the zero polynomial in column $0$.
Thus column $0$ may be replaced by an arbitrary point of $\C^I$.  Repeat
the argument for columns $1,\ldots,L-1$.  After the final step, $R$ vanishes
on all of $\C^{I\times\{0,\ldots,L-1\}}$, and hence $R=0$.
\end{proof}

\begin{lemma}[Arithmetic criterion]\label{lem:arithmetic}
Let $I$ be finite and set $q_j=\exp(-2\pi i s_j)$.  If $1$ and
$(s_j)_{j\in I}$ are linearly independent over $\Q$, then
$\alpha\mapsto q^\alpha$ is injective on $\N^I$.
\end{lemma}

\begin{proof}
Suppose $q^\alpha=q^\beta$.  Then
\[
  \exp\left(-2\pi i\sum_{j\in I}(\alpha_j-\beta_j)s_j\right)=1,
\]
so
\[
  \sum_{j\in I}(\alpha_j-\beta_j)s_j\in\Z.
\]
This is a rational linear relation among $1$ and the $s_j$.  Their linear
independence forces $\alpha_j=\beta_j$ for every $j$.
\end{proof}

\section{Simultaneous branch deformation}\label{sec:deformation}

Fix $N\geq0$.  For parameters
\[
  c=(c_{j,n})_{j\in I,\ 0\leq n\leq N}
  \in\C^{I\times\{0,\ldots,N\}},
\]
define
\begin{equation}\label{eq:deformation}
  F_{j,k}^{c}(a)=\zeta(s_j+k,a)
  +\sum_{n=0}^{N}(c_{j,n}-1)(a+n)^{-s_j-k},
  \qquad 0\leq k\leq N.
\end{equation}
The original ladder is $F^{\mathbf 1}$.

\begin{proposition}[Analyticity of the deformation]\label{prop:analyticity}
If $\operatorname{Re}s_j>1$, then every function
$a\mapsto F_{j,k}^{c}(a)$ is analytic on $\HH$.
\end{proposition}

\begin{proof}
The series \eqref{eq:hurwitz-series} is normally convergent on compact
subsets of $\HH$.  Indeed, on a compact set the finitely many initial terms
are analytic principal powers, while the tail is dominated uniformly by a
constant multiple of $m^{-\operatorname{Re}s_j-k}$.  The finite correction
in \eqref{eq:deformation} is analytic on $\HH$ because every $a+n$ lies off
the principal logarithm cut.  The conclusion follows.
\end{proof}

For a fixed branch point $-n$, define the column update
\begin{equation}\label{eq:column-update}
  (\mathcal T_n c)_{j,r}=
  \begin{cases}
    q_jc_{j,n},&r=n,\\
    c_{j,r},&r\neq n.
  \end{cases}
\end{equation}

\begin{proposition}[One-branch simultaneous monodromy]
\label{prop:one-branch}
Let
\[
  P\in\C[a][Y_{j,k}:j\in I,\ 0\leq k\leq N]
\]
and suppose that
\begin{equation}\label{eq:deformed-relation}
  P(a;F^c(a))=0\qquad(a\in\HH).
\end{equation}
Then, for every $n\in\{0,\ldots,N\}$,
\begin{equation}\label{eq:moved-relation}
  P(a;F^{\mathcal T_n c}(a))=0\qquad(a\in\HH).
\end{equation}
\end{proposition}

\begin{proof}
Fix $n$ and put
\[
  a=-n+e^z.
\]
On the upper strip $0<\operatorname{Im}z<\pi$, the $n$th singular summand
separates as
\begin{equation}\label{eq:local-model}
  F_{j,k}^{c}(-n+e^z)
  =c_{j,n}e^{-(s_j+k)z}+T_{j,k,n}^{c}(e^z).
\end{equation}
Here $T_{j,k,n}^{c}$ is holomorphic for $|e^z|<1/2$.  To see this, remove
the $n$th summand.  For $m\neq n$, the function
\[
  w\longmapsto(m-n+w)^{-s_j-k}
\]
has a holomorphic branch on $|w|<1/2$ obtained by continuation from the
upper half-disc.  The remaining infinite series converges normally there,
because $|m-n+w|$ is bounded below by a fixed positive multiple of
$1+|m-n|$ outside finitely many indices.  The finite deformation correction
has the same property.

Substituting \eqref{eq:local-model} into $P$ gives a function $G(z)$
holomorphic on the connected half-plane
\[
  \Omega=\{z\in\C:\operatorname{Re}z<-\log2\}.
\]
By \eqref{eq:deformed-relation}, it vanishes on the nonempty open subset
$\Omega\cap\{0<\operatorname{Im}z<\pi\}$.  The identity theorem therefore
gives $G\equiv0$ on $\Omega$.

Translation by $2\pi i$ fixes $e^z$, the coefficient variable
$a=-n+e^z$, and every tail term.  The singular term transforms by
\[
  e^{-(s_j+k)(z+2\pi i)}
  =e^{-2\pi i s_j}e^{-2\pi i k}e^{-(s_j+k)z}
  =q_je^{-(s_j+k)z}.
\]
Thus $G(z+2\pi i)=0$ is exactly the relation with column $n$ replaced by
the updated column in \eqref{eq:column-update}.  Returning to the $a$-plane
gives \eqref{eq:moved-relation} first on a punctured upper half-disc at
$-n$.  Proposition~\ref{prop:analyticity} and the identity theorem extend it
to all of $\HH$.
\end{proof}

\begin{corollary}[The monodromy product grid]\label{cor:grid}
If
\[
  P(a;F^{\mathbf1}(a))=0\qquad(a\in\HH),
\]
then
\begin{equation}\label{eq:grid-relation}
  P(a;F^{c(m)}(a))=0\qquad(a\in\HH)
\end{equation}
for every $m=(m_0,\ldots,m_N)\in\N^{N+1}$, where
\begin{equation}\label{eq:grid-parameters}
  c(m)_{j,n}=q_j^{m_n}.
\end{equation}
\end{corollary}

\begin{proof}
Apply Proposition~\ref{prop:one-branch} exactly $m_n$ times at the branch
point $-n$.  Updates in different columns commute because they change
disjoint coordinates.
\end{proof}

\section{Proof of ladder independence}\label{sec:main-proof}

\begin{proof}[Proof of Theorem~\ref{thm:polynomial}]
Let $P$ satisfy \eqref{eq:polynomial-relation}.  In the notation of
Section~\ref{sec:deformation}, this is
\[
  P(a;F^{\mathbf1}(a))=0\qquad(a\in\HH).
\]
Fix $a\in\HH$.  Corollary~\ref{cor:grid} shows that the polynomial
\begin{equation}\label{eq:parameter-polynomial}
  c\longmapsto P(a;F^c(a))
\end{equation}
vanishes at every product-grid point \eqref{eq:grid-parameters}.  Condition
$\mc$ and Lemma~\ref{lem:product-grid} imply that
\eqref{eq:parameter-polynomial} is the zero polynomial in all parameters
$c_{j,n}$.

It remains to calculate the range of the affine deformation map
\begin{equation}\label{eq:affine-map}
  c\longmapsto(F_{j,k}^c(a))_{j,k}.
\end{equation}
For each fixed $j$, its linear part is the $(N+1)\times(N+1)$ matrix
\begin{equation}\label{eq:W}
  W^{(j)}(a)=\bigl(W^{(j)}_{k,n}(a)\bigr)_{0\leq k,n\leq N},
  \qquad
  W^{(j)}_{k,n}(a)=(a+n)^{-s_j-k}.
\end{equation}
Set $x_n=(a+n)^{-1}$.  Then
\[
  W^{(j)}_{k,n}(a)=(a+n)^{-s_j}x_n^k,
\]
so
\begin{equation}\label{eq:det-W}
  \det W^{(j)}(a)
  =\left(\prod_{n=0}^{N}(a+n)^{-s_j}\right)
   \prod_{0\leq n<m\leq N}(x_m-x_n).
\end{equation}
Every factor in the first product is nonzero, and the $x_n$ are pairwise
distinct.  Hence \eqref{eq:det-W} is nonzero.

The complete linear part of \eqref{eq:affine-map} is block diagonal with
blocks $W^{(j)}(a)$.  It is therefore invertible, and
\eqref{eq:affine-map} is an affine automorphism of
$\C^{I\times\{0,\ldots,N\}}$.  Since \eqref{eq:parameter-polynomial}
vanishes for every $c$, it follows that
\begin{equation}\label{eq:all-y}
  P(a;y)=0
  \qquad
  \text{for every }y\in\C^{I\times\{0,\ldots,N\}}.
\end{equation}

Regard $P$ as a polynomial in the variables $Y_{j,k}$ with coefficients in
$\C[a]$.  Equation \eqref{eq:all-y} says that each coefficient, evaluated at
the chosen $a$, is zero.  Since $a$ was arbitrary in the infinite set
$\HH$, every coefficient polynomial vanishes identically.  Thus $P=0$.
\end{proof}

\begin{proposition}[Rational coefficients]\label{prop:rational}
Under condition $\mc$, the entries of $\mathcal L_N(s)$ are algebraically
independent over $\C(a)$.
\end{proposition}

\begin{proof}
Let
\[
  Q\in\C(a)[Y_{j,k}:j\in I,\ 0\leq k\leq N]
\]
be a relation.  Only finitely many coefficients of $Q$ are nonzero.  Write
each in reduced form and let $D(a)\in\C[a]$ be the product of their
denominators.  Multiplication by $D$ produces
\[
  \widetilde Q\in\C[a][Y_{j,k}:j\in I,\ 0\leq k\leq N].
\]
Because $\C[a]$ is an integral domain, $\widetilde Q\neq0$ whenever
$Q\neq0$.  The relation defined by $Q$ gives the relation defined by
$\widetilde Q$ away from the finite zero set of $D$.  Its left-hand side is
analytic on $\HH$, so the identity theorem extends the cleared relation
across those points.  Theorem~\ref{thm:polynomial} gives
$\widetilde Q=0$, hence $Q=0$.
\end{proof}

\section{Differential and arbitrary-family consequences}
\label{sec:differential}

\begin{lemma}[Jet-ladder identity]\label{lem:jet-ladder}
If $\operatorname{Re}s>1$, then for every $k\in\N$ and $a\in\HH$,
\begin{equation}\label{eq:jet-ladder}
  \partial_a^k\zeta(s,a)=(-1)^k(s)_k\zeta(s+k,a).
\end{equation}
Moreover, $(s)_k\neq0$.
\end{lemma}

\begin{proof}
Normal convergence of \eqref{eq:hurwitz-series} on compact subsets of
$\HH$ permits termwise differentiation.  One differentiation gives
\[
  \partial_a\zeta(s,a)=-s\zeta(s+1,a).
\]
Induction gives \eqref{eq:jet-ladder}.  If $(s)_k=0$, then
$s\in\{0,-1,\ldots,-k+1\}$, which contradicts
$\operatorname{Re}s>1$.
\end{proof}

\begin{proof}[Proof of Theorem~\ref{thm:main}]
Proposition~\ref{prop:rational} proves algebraic independence of every finite
ladder.  By Lemma~\ref{lem:jet-ladder}, a finite jet rectangle
\[
  (\partial_a^k\zeta(s_j,a))_{j\in I,\ 0\leq k\leq N}
\]
is obtained from $\mathcal L_N(s)$ by an invertible diagonal coordinate
change over $\C$.  It is therefore algebraically independent over $\C(a)$.

To express the conclusion inside $\Minf$, suppose a polynomial over $\C(a)$
in finitely many differential-field jets vanishes in $\Minf$.  Clear its
coefficient denominators.  The resulting analytic upper-end germ vanishes on
some $\HH_H$.  The identity theorem extends this vanishing to $\HH$.
Theorem~\ref{thm:polynomial}, followed by the invertible diagonal change from
Lemma~\ref{lem:jet-ladder}, forces the polynomial to be zero.  Hence the
Hurwitz germs are differentially algebraically independent over $\C(a)$.
\end{proof}

\begin{proof}[Proof of Corollary~\ref{cor:arithmetic}]
Apply Lemma~\ref{lem:arithmetic} and Theorem~\ref{thm:main}.
\end{proof}

\begin{proof}[Proof of Corollary~\ref{cor:arbitrary}]
Every polynomial relation involves finitely many indices
$J\subset K$ and finitely many derivative orders.  Rational linear
independence of the augmented family $1,(s_j)_{j\in K}$ passes to the
subfamily $1,(s_j)_{j\in J}$.  Corollary~\ref{cor:arithmetic} kills the
restricted relation.  This proves the jet statement.  The ladder statement
follows from the same finite restriction and Lemma~\ref{lem:jet-ladder}.
\end{proof}

\section{The resonance boundary}\label{sec:boundary}

Condition $\mc$ has an exact orbit-theoretic meaning.  If it fails, there
exist $\alpha\neq\beta$ in $\N^I$ such that
\[
  q^\alpha=q^\beta.
\]
The nonzero polynomial $X^\alpha-X^\beta$ then vanishes on the common cyclic
orbit.  Thus the orbit lies in a proper algebraic subvariety of
$\C^I$, and the product-grid argument supplies only the corresponding
smaller deformation variety at each branch point.

This observation is a boundary of the theorem, not a converse.  Failure of
$\mc$ does not by itself produce a differential relation among the Hurwitz
functions.  It says that cyclic monodromy alone no longer spans the full
parameter space.  The remaining positive datum is the resonance lattice
\[
  \left\{u\in\Z^I:\sum_{j\in I}u_js_j\in\Z\right\},
\]
which determines the character relations visible to the common loop.  A
classification in resonant chambers would require additional analytic
invariants beyond the cyclic power orbit.

The restriction $\operatorname{Re}s_j>1$ is used to work directly with the
normally convergent Hurwitz series and to ensure that all Pochhammer factors
in the jet-ladder change of coordinates are nonzero.  No assertion is made
here about continuation of the theorem to exponents outside that half-plane.

\section{Conclusion}

The joint differential algebra of fixed-exponent Hurwitz zeta functions in
the shift variable is controlled by a simple monodromy-character condition.
At one branch point, the exponent coordinates share one loop count.  Character
separation makes that constrained orbit Zariski dense.  Across different
branch points, the loop counts are independent, and a block-Vandermonde
automorphism converts the resulting product geometry into full finite-jet
controllability.  This yields differential algebraic independence over
$\C(a)$ for every finite nonresonant family and, under rational linear
independence, for arbitrary index families in the finite-subfamily sense.

\begin{thebibliography}{99}

\bibitem{ADH}
B.~Adamczewski, T.~Dreyfus, and C.~Hardouin,
\emph{Hypertranscendence and linear difference equations},
J. Amer. Math. Soc. \textbf{34} (2021), 475--503.
\href{https://doi.org/10.1090/jams/960}{doi:10.1090/jams/960}.

\bibitem{AmouKatsurada}
M.~Amou and M.~Katsurada,
\emph{Differential transcendence of a class of generalized Dirichlet series},
Illinois J. Math. \textbf{45} (2001), 939--948.
\href{https://doi.org/10.1215/ijm/1258138161}{doi:10.1215/ijm/1258138161}.

\bibitem{CostinGaroufalidis}
O.~Costin and S.~Garoufalidis,
\emph{Resurgence of the fractional polylogarithms},
Math. Res. Lett. \textbf{16} (2009), 817--826.
\href{https://doi.org/10.4310/MRL.2009.v16.n5.a5}{doi:10.4310/MRL.2009.v16.n5.a5}.

\bibitem{DLMF}
NIST Digital Library of Mathematical Functions,
\emph{Hurwitz zeta function}, Section 25.11.
\url{https://dlmf.nist.gov/25.11}.

\bibitem{Joyner}
S.~Joyner,
\emph{On an algebraic version of the Knizhnik-Zamolodchikov equation},
arXiv:1009.5659 (2010).
\url{https://arxiv.org/abs/1009.5659}.

\bibitem{LagariasLiII}
J.~C. Lagarias and W.-C.~W. Li,
\emph{The Lerch zeta function II. Analytic continuation},
Forum Math. \textbf{24} (2012), 49--84.
\href{https://doi.org/10.1515/FORM.2011.048}{doi:10.1515/FORM.2011.048}.

\bibitem{LagariasLiIII}
J.~C. Lagarias and W.-C.~W. Li,
\emph{The Lerch zeta function III. Polylogarithms and special values},
Res. Math. Sci. \textbf{3} (2016), Article 2.
\url{https://arxiv.org/abs/1506.06161}.

\bibitem{Laurincikas}
A.~Laurin\v{c}ikas,
\emph{Functional independence of periodic Hurwitz zeta functions},
Math. Notes \textbf{83} (2008), 65--71.
\href{https://doi.org/10.1134/S0001434608010082}{doi:10.1134/S0001434608010082}.

\bibitem{Ostrowski}
A.~Ostrowski,
\emph{\"Uber Dirichletsche Reihen und algebraische Differentialgleichungen},
Math. Z. \textbf{8} (1920), 241--298; English translation by E.~C. Hansen,
Y.~Stone, and J.~Wolfson, arXiv:2211.02088.
\url{https://arxiv.org/abs/2211.02088}.

\end{thebibliography}

\end{document}
